problem:  
Let \(M^m \hookrightarrow (\mathbb{S}^n,g_{\mathrm{std}})\) be a compact minimal submanifold with Jacobi operator \(J_M\) acting on normal vector fields.  
For a smooth mean‑zero function \(u\) on \(\mathbb{S}^n\), consider the infinitesimal conformal deformation of the ambient metric \(g_t = e^{2t u}g_{\mathrm{std}}\) and denote by \(J_{M,t}\) the corresponding family of Jacobi operators.  The eigenvalues of \(J_{M,t}\) form continuous curves \(\lambda_i(t)\) with \(\lambda_i(0)\) being the spectrum of \(J_M\).  

**Problem:**  
Assume that \(\lambda_1(0)=0\) is a *simple zero‑eigenvalue* of \(J_M\).  Determine the generic sign and asymptotic behavior of the first derivative \(\dot \lambda_1(0)\) under all mean‑zero conformal directions \(u\).  
In particular, does *every* simple zero of the Jacobi operator at a minimal immersion necessarily become positive (stable) or negative (unstable) under all small conformal perturbations of the ambient metric, or is the sign of \(\dot\lambda_1(0)\) genuinely dependent on intrinsic conformal data of \(u\) and the geometry of \(M\)?  

Further, investigate whether there are geometric or spectral conditions on \(M\) guaranteeing that   
\[
\dot\lambda_1(0) = 0 \quad \text{for all mean‑zero } u,
\]
so that the spectrum at zero remains *conformally stationary*.  

Why is it a "good" problem:  
This reformulation replaces the ill‑defined “index growth rate” by the precise analytical notion of how eigenvalues of the Jacobi operator move under an infinitesimal conformal change of the ambient metric.  It explores the *spectral flow* of stability operators for minimal immersions when the background geometry itself varies conformally.  

Resolving the question requires new understanding of the coupling between:
- the conformal deformation tensor of the sphere,
- the second variation formulas for minimal submanifolds, and
- perturbation theory for eigenvalues of geometric elliptic operators.

The problem stands at the junction of minimal submanifold theory, conformal geometry, and spectral analysis.  Its statement is simple—“What is the direction of motion of zero eigenvalues of the Jacobi operator under conformal metric deformations?”—but the answer would reveal whether conformal geometry inherently stabilizes or destabilizes minimal embeddings.  This depth, clarity, and interdisciplinary reach make it a strong and insightful research problem.